\documentclass[12pt]{article}

\setlength{\parindent}{0pt}

\usepackage{amssymb, amsmath, fullpage, graphicx}

\begin{document}

\title{\bf{Aufgaben zum Vortrag}}
\author{Jan Wiesemann}
\date{Wintersemester 2024/25}
\maketitle

\subsection*{Aufgabe 1}
Prüfe im Detail nach, dass die im Vortrag erwähnte Abbildung
\[\varphi: \mathbb C\to\mathbb C; \hspace{3em} x+iy \longmapsto x-iy\]
wirklich einen Ringautomorphismus definiert.

\subsection*{Aufgabe 2}
Beweise oder widerlege: Enthält ein Ideal $\mathfrak a$ eine Einheit $a\in R^\times$, so ist $\mathfrak a=(1)$.

\subsection*{Aufgabe 3}
Sei $R$ ein Körper. Beweise, dass die Charakteristik $p$ von $R$ entweder 0 oder prim.

\vspace{2.5em}
QR-Code zum Skript:
\vspace{-1.5em}
\center\includegraphics{qr.pdf}

\raggedright
Falls der QR-Code nicht funktioniert:
https://github.com/janmagtoast/Ringhomomorphismen-und-Ideale
\end{document}
